\begin{center}
    {\LARGE \textbf{2020无机化学}} \\[2ex]
    {\large 时量180分钟 \quad 满分100分}
\end{center}

\vspace{0.5cm}
\noindent\textbf{注意事项:}
\begin{enumerate}[label=\arabic*., parsep=0pt, itemsep=0pt]
    \item 答题前,考生先将自己的姓名、学号、考场号、座位号填写在答题卡上,将条形码准确粘贴在考生信息条形码粘贴区。
    \item 考试结束后,将本试卷与答题纸一并收回。
    \item 回答各个问题时,将答案写在答题纸上,写在本试卷上无效。
\end{enumerate}

\vspace{0.5cm}
\noindent\textbf{一、选择题（20 pts.）}

\begin{enumerate}[label=\arabic*., itemsep=1.5ex]
    \item 下列配合物中，具有平面正方形结构的是
    \begin{tasks}(4)
        \task \ce{Fe(CO)5}
        \task \ce{[Zn(NH3)4]^{2+}}
        \task \ce{Ni(CO)4}
        \task \ce{[PtCl4]^{2-}}
    \end{tasks}
    \item 在酸性介质中，物质都不能发生歧化反应的是
    \begin{tasks}(4)
        \task \ce{SO4^{2-}}, \ce{H2SO3}
        \task \ce{H2SO3}, \ce{S2O8^{2-}}
        \task \ce{S2O3^{2-}}, \ce{S}
        \task \ce{SO4^{2-}}, \ce{S}
    \end{tasks}
    \item 下列配合物中，没有旋光异构体的是
    \begin{tasks}(4)
        \task \ce{K3[Fe(C2O4)3]}
        \task \ce{[Ni(en)3]Cl2}
        \task \ce{[Co(en)2(NH3)2]Cl2}
        \task \ce{[PtCl3(NH3)(OH)]}
    \end{tasks}
    \item 下列含氧酸盐中，溶解度最小的是
    \begin{tasks}(4)
        \task \ce{KClO}
        \task \ce{KClO2}
        \task \ce{KClO3}
        \task \ce{KClO4}
    \end{tasks}
    \item 下列各组硫化物中，颜色基本相同的是
    \begin{tasks}(4)
        \task \ce{ZnS}, \ce{MnS}, \ce{FeS}
        \task \ce{CdS}, \ce{SnS2}, \ce{As2S3}
        \task \ce{Ag2S}, \ce{Sb2S5}, \ce{Bi2S3}
        \task \ce{SnS}, \ce{PbS}, \ce{CoS}
    \end{tasks}
    \item 下列溶液中与硝酸银溶液作用生成的沉淀颜色最浅的是
    \begin{tasks}(4)
        \task \ce{H3PO3}
        \task \ce{H3PO2}
        \task \ce{HPO3}
        \task \ce{PH3}
    \end{tasks}
    \item 下列氯化物中最不稳定的是
    \begin{tasks}(4)
        \task \ce{SnCl4}
        \task \ce{SbCl5}
        \task \ce{GeCl4}
        \task \ce{PbCl4}
    \end{tasks}
    \item 在乙醇中溶解度最大的是
    \begin{tasks}(4)
        \task \ce{FeCl2}
        \task \ce{FeCl3}
        \task \ce{CoCl2}
        \task \ce{NiCl2}
    \end{tasks}
    \item 第五周期元素原子中，未成对电子数最多的为
    \begin{tasks}(4)
        \task 4个
        \task 5个
        \task 6个
        \task 7个
    \end{tasks}
    \item 下列分子中，不存在$d$-$p\pi$配键的是
    \begin{tasks}(4)
        \task \ce{SOCl2}
        \task \ce{HNO3}
        \task \ce{H3PO4}
        \task \ce{HClO2}
    \end{tasks}
\end{enumerate}

\vspace{0.5cm}
\noindent\textbf{二、写出实验现象并给出相关反应方程式（40 pts.）}

\begin{enumerate}[label=\arabic*., start=1, itemsep=1.5ex]
    \item 向\ce{KI}溶液中加入\ce{H2O2}溶液。
    \item 分别向\ce{Na4P2O7}，\ce{Na2HPO4}，\ce{NaH2PO2}溶液中加入\ce{AgNO3}溶液。
    \item 向\ce{CoCl2}溶液中滴加\ce{NaOH}溶液和碘水，充分反应后再加入硫酸。
    \item 向\ce{K2MnO4}溶液中缓慢通入\ce{SO2}。
    \item 向\ce{Hg2(NO3)2}溶液中缓慢滴加\ce{KI}溶液。
    \item 金属铬溶于盐酸。
    \item 向三氯化铁溶液中加入硫氰化钾，再加入适量二氯化锡。
    \item 将$3\,\mathrm{mol\cdot dm^{-3}}$ \ce{CoCl2}溶液加热，再滴加硝酸银溶液。
    \item 在煤气灯上加热含有结晶水的硝酸铜晶体。
    \item 氯水滴入\ce{KBr}和\ce{KI}混合溶液中。
\end{enumerate}

\vspace{0.5cm}
\noindent\textbf{三、简答题（40 pts.）}

\begin{enumerate}[label=\arabic*., start=1, itemsep=1.5ex]
    \item 为什么\ce{[Fe(CN)6]^{4-}}不如\ce{[Fe(CN)6]^{3-}}稳定？
    \item 通过计算说明\ce{Cu^+}在氨水中能否稳定存在？
    已知：
    
    $E^\ominus(\ce{Cu^{2+}/Cu}) = 0.34\,\mathrm{V}$，$E^\ominus(\ce{Cu^+/Cu}) = 0.52\,\mathrm{V}$，

    $\lg K^\ominus_{\text{稳}}(\ce{[Cu(NH3)4]^2+}) = 13.32$，$\lg K^\ominus_{\text{稳}}(\ce{[Cu(NH3)2]^+}) = 10.86$
    \item 解释下列实验事实：
    \begin{enumerate}[label=(\arabic*), itemsep=1ex]
        \item 硅不溶于硝酸而铅溶于硝酸。
        \item 铬酸铅既溶于硝酸也溶于氢氧化钠溶液。
    \end{enumerate}
    \item 请解释实验现象：向\ce{CuSO4}溶液中加入\ce{KI}溶液，有黄色沉淀生成；向黄色沉淀中加入过量\ce{Na2SO3}溶液，沉淀转化为白色；向黄色沉淀中加入过量\ce{Na2S2O3}溶液，沉淀消失，得到无色溶液。
    \item 向\ce{Na2SO3}的淀粉溶液中滴加酸性\ce{KIO3}溶液和向酸性\ce{KIO3}的淀粉溶液中滴加\ce{Na2SO3}溶液观察到的现象是否相同？为什么？给出相关的反应方程式。
    \item 为什么浓高氯酸可以氧化碘单质而稀高氯酸却不能氧化碘单质？
    \item 为什么铜与氰化钠溶液作用放出氢气，而在有氧气存在时可溶于浓氨水？
    \item 论证下列离子能否共存？写出反应方程式。
    \[
    \ce{K+}, \ce{MnO4^-}, \ce{SO4^{2-}}, \ce{Mn^{2+}}, \ce{I^-}
    \]
\end{enumerate}

\vspace{0.5cm}
\noindent\textbf{四、制备与分离（20 pts.）}

\begin{enumerate}[label=\arabic*., start=1, itemsep=1.5ex]
    \item 用反应方程式表示化学方法制取\ce{F2}的过程。
    \item 以方铅矿为原料制备二氧化铅。
    \item 不用硫化氢和硫化物，设计方案分离含有\ce{Al^{3+}}、\ce{Co^{2+}}、\ce{Cr^{3+}}、\ce{Ni^{2+}}、\ce{Fe^{3+}}的混合溶液。
    \item 为什么黄血盐可以由硫酸亚铁溶液和氰化钾溶液混合直接制备，而赤血盐却不能由三氯化铁溶液和氰化钾直接混合来制备？如何制备赤血盐？
\end{enumerate}

\vspace{0.5cm}
\noindent\textbf{五、鉴别与推断（30 pts.）}

\begin{enumerate}[label=\arabic*., start=1, itemsep=1.5ex]
    \item 选用一种试剂将下列溶液鉴别开，并进一步加以确证。
    \[
    \ce{Na2SO4}, \ce{Na2S2O3}, \ce{Na2SO3}, \ce{Na2S2}, \ce{Na2S}
    \]
    \item 用三种方法鉴别下列各对物质
    \begin{enumerate}[label=(\arabic*), itemsep=1ex]
        \item \ce{SbCl3}和\ce{Bi(NO3)3}
        \item \ce{MnO2}和\ce{PbO2}
    \end{enumerate}
    \item 请回答下列问题：
    \begin{enumerate}[label=(\arabic*), itemsep=1ex]
        \item 在\ce{Pt}制的坩埚底部和壁上粘有高温灼烧过的\ce{Al2O3}，应如何清洗之？
        \item 提纯含有少量铁、钴杂质的金属镍。
        \item 提纯含有少量\ce{Ni^{2+}}杂质的\ce{Co^{2+}}溶液
    \end{enumerate}
    \item 无色晶体\ce{A}溶于稀盐酸得无色溶液，再加入氢氧化钠溶液得白色沉淀\ce{B}，\ce{B}溶于过量氢氧化钠溶液得无色溶液\ce{C}。将\ce{B}溶于盐酸后蒸发、浓缩后又析出\ce{A}。向\ce{A}的稀盐酸溶液加入硫化氢溶液生成橙色沉淀\ce{D}。\ce{D}与过硫化钠可发生氧化还原反应，反应物之间相互作用得到无色溶液\ce{E}。将\ce{A}置于水中生成白色沉淀\ce{F}。写出\ce{A-F}所代表的物质的化学式，并用化学方程式表示$B\rightarrow C$和$D\rightarrow E$的过程。
    \item 黑蓝色化合物\ce{A}不溶于水，但可溶于稀盐酸生成\ce{B}的溶液，稀释该溶液时，析出白色沉淀\ce{C}，酸化时，沉淀溶解。向\ce{B}的溶液中通硫化氢时，生成灰黑色沉淀\ce{D}，该沉淀可溶于浓盐酸。\ce{D}与硝酸反应时产生浅黄色沉淀\ce{E}，无色气体\ce{F}和化合物\ce{G}的溶液。\ce{F}在空气中迅速变成红棕色。向\ce{G}的溶液中通硫化氢时，生成黄色沉淀\ce{H}。试写出\ce{A}，\ce{B}，\ce{C}，\ce{D}，\ce{G}和\ce{H}所代表的物质的化学式，并写出\ce{D}与硝酸反应的化学方程式。
    \item 将钠盐\ce{A}溶液滴加到\ce{B}溶液（硝酸盐）中，立即有白色沉淀\ce{C}生成，\ce{C}放置一段时间会变为黑色沉淀\ce{D}；若将\ce{B}溶液滴加到\ce{A}溶液中，开始形成的是无色溶液\ce{E}；当加入相当多\ce{B}后才会有\ce{C}沉淀出来。向溶液\ce{B}中加入钾盐\ce{F}溶液，可以生成浅黄色沉淀\ce{G}，\ce{G}可溶于\ce{A}溶液中，\ce{G}在光照下可分解变黑。若向溶液\ce{E}中加入盐酸，有大量沉淀\ce{H}生成，放出刺激性气体\ce{I}。沉淀\ce{H}溶于过量浓硝酸。试写出\ce{A} $\sim$ \ce{I}的化学式。
\end{enumerate}